\documentclass[11pt,a4paper]{article}
\usepackage[utf8]{inputenc}
\usepackage[margin=1in]{geometry}
\usepackage{hyperref}
\usepackage{enumitem}
\usepackage{titlesec}
\usepackage{xcolor}
\usepackage{wasysym}

\definecolor{sbu-red}{RGB}{153, 0, 0}
\definecolor{slate-grey}{RGB}{70, 80, 95}

\hypersetup{
    colorlinks=true,
    linkcolor=sbu-red,
    filecolor=sbu-red,      
    urlcolor=sbu-red,
    pdftitle={Summer Systems Reading List},
}

\titleformat{\section}{\large\bfseries\color{sbu-red}}{}{0em}{}[\titlerule]
\titleformat{\subsection}{\normalsize\bfseries\color{slate-grey}}{}{0em}{}

\title{\textbf{Summer Reading List: Operating Systems \& File Systems}}
\author{\textbf{Stony Brook University} \\ Graduate Systems Reading Program}
\date{Summer 2026}

\begin{document}

\maketitle

\begin{abstract}
This curated bibliography compiles 26 seminal and modern academic papers covering foundational operating systems, innovative kernel architectures, local and distributed file systems, virtualization, and hardware-software co-design. It includes classic texts essential for graduate studies alongside prominent research originating from Stony Brook University labs.
\end{abstract}

\section{Part I: Foundational \& Monolithic OS Architectures}
These classic papers established fundamental abstractions such as processes, threads, virtual memory, and unified resource interfaces.

\begin{enumerate}
    \item \textbf{The UNIX Time-Sharing System} \\
    \textit{Dennis M. Ritchie and Ken Thompson (1974)} \\
    The foundational text for UNIX. Introduces the paradigm of treating "everything as a file," hierarchical file systems, and shell pipelines. \\
    $\CheckedBox$
    
    \item \textbf{The Nucleus of a Multiprogramming System} \\
    \textit{Per Brinch Hansen (1970)} \\
    An early, highly influential exploration of separating operating system policy from mechanism, serving as a philosophical precursor to microkernels. \\
    $\Box$
    
    \item \textbf{The Protection of Information in Computer Systems} \\
    \textit{Jerome H. Saltzer and Michael D. Schroeder (1975)} \\
    The definitive security architecture paper. Establishes timeless design principles including "least privilege" and "fail-safe defaults." \\
    $\Box$
    
    \item \textbf{Hints for Computer System Design} \\
    \textit{Butler W. Lampson (1983)} \\
    A distillation of practical wisdom for system builders, championing principles like "keep it simple," "use hints," and "plan to throw one away." \\
    $\Box$
\end{enumerate}

\section{Part II: Kernel Design Evolutions (Microkernels \& Extensibility)}
Key milestones in the historical structural debate regarding whether the kernel should be monolithic or a minimal IPC broker.

\begin{enumerate}[resume]
    \item \textbf{The Performance of Micro-Kernel Implementations} \\
    \textit{Jochen Liedtke (1993)} \\
    The landmark paper that revitalized microkernel research by introducing L4 and demonstrating that IPC bottlenecks were implementation flaws, not architectural limitations. \\
    $\Box$
    
    \item \textbf{Extensibility, Safety and Performance in the SPIN Operating System} \\
    \textit{Brian N. Bershad, Stefan Savage, Pardha Pratim, Emin G\"{u}n Sirer, Marc E. Fiuczynski, David Becker, Craig Chambers, and Susan Eggers (1995)} \\
    Explored dynamically downloading user code safely into kernel space via type-safe languages (Modula-3), serving as a spiritual ancestor to eBPF. \\
    $\Box$
    
    \item \textbf{Exokernel: An Operating System Architecture for Application-Specific Resource Management} \\
    \textit{Dawson R. Engler, M. Frans Kaashoek, and James O'Toole Jr. (1995)} \\
    Advocated for eliminating standard abstractions to give untrusted applications low-level, direct access to physical hardware resources. \\
    $\Box$
\end{enumerate}

\section{Part III: Classical and Local File Systems}
The evolution of persistent storage mapping physical layouts to handle physical mechanical constraints and data corruption.

\begin{enumerate}[resume]
    \item \textbf{A Fast File System for UNIX} \\
    \textit{Marshall Kirk McKusick, William N. Joy, Samuel J. Leffler, and Robert S. Fabry (1984)} \\
    The classic BSD FFS paper detailing how optimizing block layout into "cylinder groups" drastically improved physical disk locality and bandwidth. \\
    $\Box$
    
    \item \textbf{The Design and Implementation of a Log-Structured File System} \\
    \textit{Mendel Rosenblum and John K. Ousterhout (1992)} \\
    Introduced LFS, treating disk storage as a continuous, append-only log to optimize write throughput; highly applicable to modern Flash Translation Layers (FTL). \\
    $\Box$
    
    \item \textbf{Analysis and Evolution of Journaling File Systems} \\
    \textit{Vijayan Prabhakaran, Andrea C. Arpaci-Dusseau, and Remzi H. Arpaci-Dusseau (2005)} \\
    An exhaustive analysis of transactional logging variants (data vs. metadata journaling) under failure states to preserve consistency. \\
    $\Box$
\end{enumerate}

\section{Part IV: Distributed Systems \& Network File Systems}
Architectures adapting state, caching, consistency protocols, and scale across network boundaries.

\begin{enumerate}[resume]
    \item \textbf{Scale and Performance in a Distributed File System} \\
    \textit{John H. Howard, Michael J. Kazar, Sherri G. Menees, David A. Nichols, Mahadev Satyanarayanan, Robert N. Sidebotham, and Michael J. West (1988)} \\
    Introduced the Andrew File System (AFS), detailing aggressive client-side whole-file caching and stateful callback architectures. \\
    $\Box$
    
    \item \textbf{The Google File System} \\
    \textit{Sanjay Ghemawat, Howard Gobioff, and Shun-Tak Leung (2003)} \\
    Pioneered warehouse-scale storage by assuming commodity hardware failure is the norm, optimizing for multi-gigabyte, append-only workloads. \\
    $\Box$
    
    \item \textbf{Bigtable: A Distributed Storage System for Structured Data} \\
    \textit{Fay Chang, Jeffrey Dean, Sanjay Ghemawat, Wilson C. Hsieh, Deborah A. Wallach, Mike Burrows, Tushar Chandra, Andrew Fikes, and Robert E. Gruber (2006)} \\
    A wide-column, highly scalable database layer built on top of GFS using Log-Structured Merge-trees (LSM-trees). \\
    $\Box$
    
    \item \textbf{Ceph: A Scalable, High-Performance Distributed File System} \\
    \textit{Sage A. Weil, Scott A. Brandt, Ethan L. Miller, Darrell D. E. Long, and Carlos Maltzahn (2006)} \\
    Decoupled storage metadata maps via the algorithmic CRUSH pseudo-random data placement function, removing structural centralized lookup tables. \\
    $\Box$
\end{enumerate}

\section{Part V: Virtualization, Concurrency, \& Modern Paradigms}
Adapting operating systems to massive core scaling, direct memory buses, and single-tenant cloud execution contexts.

\begin{enumerate}[resume]
    \item \textbf{SOSP's Do's and Don'ts: The Virtual Machine Monitor} \\
    \textit{Mendel Rosenblum and Tal Garfinkel (2005)} \\
    A retrospective and forward-looking text outlining the practical industrial rebirth of full virtualization and hypervisors. \\
    $\Box$
    
    \item \textbf{Xen and the Art of Virtualization} \\
    \textit{Paul Barham, Boris Dragovic, Keir Fraser, Steven Hand, Tim Harris, Alex Ho, Rolf Neugebauer, Ian Pratt, and Andrew Warfield (2003)} \\
    Introduced paravirtualization, requiring hypervisor-aware modifications to guest OS kernels to hit near-native computation performance. \\
    $\Box$
    
    \item \textbf{The Multikernel: A New OS Architecture for Scalable Multicore Systems} \\
    \textit{Andrew Baumann, Paul Barham, Pierre-Evariste Dagand, Tim Harris, Rebecca Isaacs, Simon Peter, Timothy Roscoe, Adrian Sch\"{u}pbach, and Akhilesh Singhania (2009)} \\
    The Barrelfish paper, treating cache-coherent multicore chips as a network of message-passing distributed nodes rather than a unified global shared memory pool. \\
    $\Box$
    
    \item \textbf{Non-Volatile Memory} \\
    \textit{Steven Pelley, Thomas F. Wenisch, Brian T. Gold, and David Bridge (2014)} \\
    Analyzes the architectural impacts of placing fast persistent memory directly on the processor's memory bus, collapsing storage and volatile caches. \\
    $\Box$
    
    \item \textbf{OSV---An Operating System for the Cloud} \\
    \textit{Avi Kivity, Dor Laor, Glauber Costa, Pekka Enberg, Nadav Har'El, Don Marti, and Vlad Zolotarov (2014)} \\
    Pioneered modern cloud "unikernels" by removing ancient multi-user isolation constructs to streamline runtime environments inside virtualized infrastructure. \\
    $\Box$
    
    \item \textbf{Verifying a File System Implementation} \\
    \textit{Konstantine Arkoudas, Karen Zee, Viktor Kuncak, and Martin Rinard (2004)} \\
    Presents rigorous formal systems verification methods using simulation relations to mathematically validate structural correctness against data corruption. \\
    $\Box$
\end{enumerate}

\section{Part VI: Selected Systems Research from Stony Brook University}
Advanced research originating from labs at Stony Brook University (COMPAS, FSL, and associated collaborations), bridging hardware accelerators, modern microarchitecture, and file systems.

\begin{enumerate}[resume]
    \item \textbf{Clearing the Clouds: A Study of Emerging Scale-out Workloads} \\
    \textit{Michael Ferdman, Almutaz Adileh, Onur Kocatyali, Volkan Volkan, Jevgenijs Jaleckis, Djordje Jevdjic, Saurabh Kestur, Per Ekman, Adrianne Castro, Mark D. Hill, and Babak Falsafi (2012)} \\
    Winner of the ASPLOS Test-of-Time Award. Highlights the disconnect between conventional CPU architectures and modern data-center scales. \\
    $\Box$
    
    \item \textbf{A Case for Hardware-Managed Virtual Memory on FPGAs} \\
    \textit{Zeljko Korolija, Timothy Roscoe, and Michael Ferdman (2020)} \\
    Investigates decoupling hardware accelerators from core kernel virtualization mechanics by placing memory management units directly inside FPGAs. \\
    $\Box$
    
    \item \textbf{TailCheck: A Lightweight Heap Overflow Detection Mechanism} \\
    \textit{Seonmyeong So, Gyuwan Kim, and Michael Ferdman (2022)} \\
    Demonstrates hardware-software co-design using specialized structural tags to eliminate heap memory safety threats with negligible runtime overhead. \\
    $\Box$
    
    \item \textbf{To FUSE or Not to FUSE: Performance of User-Space File Systems} \\
    \textit{Bharath Kumar Reddy Vangoor, Vasily Tarasov, and Erez Zadok (2017)} \\
    Provides a deep structural audit of performance degradation patterns and memory copying boundaries in standard user-space file systems layers. \\
    $\Box$
    
    \item \textbf{Performance and Resource Utilization of FUSE User-Space File Systems} \\
    \textit{Vasily Tarasov, Bharath Kumar Reddy Vangoor, and Erez Zadok (2019)} \\
    Extends user-space layout profiling to show specific interaction mechanics during high context-switching and concurrent multi-threaded I/O workloads. \\
    $\Box$
    
    \item \textbf{Evaluating Performance and Energy in File System Server Workloads} \\
    \textit{Priya Sehgal, Ming Chen, and Erez Zadok (2010)} \\
    Examines the tangible electrical footprints and efficiency boundaries of competing Linux storage layouts across intensive cloud servers. \\
    $\Box$
\end{enumerate}

\end{document}
